% LaTeX figures for different factuality evaluation approach prompts

\documentclass[11pt]{article}
\usepackage{listings}
\usepackage{xcolor}
\usepackage{tcolorbox}
\usepackage{geometry}
\usepackage{caption}
\usepackage{subcaption}

\geometry{margin=1in}

% Define colors for prompt formatting
\definecolor{promptbg}{RGB}{248,249,250}
\definecolor{promptborder}{RGB}{200,200,200}
\definecolor{responsecolor}{RGB}{0,123,255}
\definecolor{systemcolor}{RGB}{108,117,125}

% Create a nice prompt box style
\tcbuselibrary{listings,skins,breakable}

\newtcblisting{promptbox}[1][]{
  listing only,
  boxrule=0.5pt,
  colback=promptbg,
  colframe=promptborder,
  sharp corners,
  left=5mm,
  right=5mm,
  top=3mm,
  bottom=3mm,
  breakable,
  listing options={
    basicstyle=\small\ttfamily,
    breaklines=true,
    columns=flexible,
    keepspaces=true,
    showstringspaces=false,
    commentstyle=\color{systemcolor}\itshape,
    keywordstyle=\color{responsecolor}\bfseries,
  },
  #1
}

\begin{document}

% ============================================================================
% BILATERAL EVALUATION PROMPTS
% ============================================================================

\begin{figure}[htbp]
\centering
\begin{minipage}{0.48\textwidth}
\begin{promptbox}[title={Verification Prompt (P⁺)}]
You are tasked with determining whether the following assertion can be verified as true based on available evidence and knowledge.

Assertion: {assertion}

Your task is to determine if this assertion can be verified. Consider all available evidence, facts, and reliable sources of information.

You must respond with exactly one of these two token sequences:
- VERIFIED (if the assertion can be confirmed as true based on evidence)
- CANNOT VERIFY (if the assertion cannot be confirmed as true, either due to lack of evidence, uncertainty, or because it is false)

Do not provide any explanation or additional text. Respond with only the required token sequence.

Response:
\end{promptbox}
\end{minipage}
\hfill
\begin{minipage}{0.48\textwidth}
\begin{promptbox}[title={Refutation Prompt (P⁻)}]
You are tasked with determining whether the following assertion can be refuted (shown to be false) based on available evidence and knowledge.

Assertion: {assertion}

Your task is to determine if this assertion can be refuted. Consider all available evidence, facts, and reliable sources of information that might contradict the assertion.

You must respond with exactly one of these two token sequences:
- REFUTED (if the assertion can be shown to be false based on evidence)
- CANNOT REFUTE (if the assertion cannot be shown to be false, either due to lack of contradictory evidence, uncertainty, or because it is true)

Do not provide any explanation or additional text. Respond with only the required token sequence.

Response:
\end{promptbox}
\end{minipage}
\caption{Bilateral Evaluation Prompts. The bilateral approach uses two separate prompts to independently assess verifiability (P⁺) and refutability (P⁻) of assertions, resulting in generalized truth values ⟨u,v⟩.}
\label{fig:bilateral-prompts}
\end{figure}

% ============================================================================
% UNILATERAL EVALUATION PROMPTS
% ============================================================================

\begin{figure}[htbp]
\centering
\begin{minipage}{0.48\textwidth}
\begin{promptbox}[title={Unilateral-Binary (Forced Choice)}]
Determine whether the following statement is correct.
Conclude with a single line containing ONLY one of these two words:
CORRECT
INCORRECT

Statement: {assertion}
\end{promptbox}
\end{minipage}
\hfill
\begin{minipage}{0.48\textwidth}
\begin{promptbox}[title={Unilateral-Ternary (With Uncertainty)}]
Determine whether the following statement is correct.

You have three options:
1. CORRECT - if you are confident the statement is true
2. INCORRECT - if you are confident the statement is false  
3. UNCERTAIN - if you lack sufficient confidence to determine correctness

Important: Only respond CORRECT or INCORRECT if you are reasonably confident.
If you have doubts or the statement is ambiguous, respond UNCERTAIN.

Conclude with a single line containing ONLY one of these three words:
CORRECT
INCORRECT
UNCERTAIN

Statement: {assertion}
\end{promptbox}
\end{minipage}
\caption{Unilateral Evaluation Prompts. Unilateral-Binary forces a binary choice with 100\% coverage, while Unilateral-Ternary allows abstention through an UNCERTAIN response.}
\label{fig:unilateral-prompts}
\end{figure}

% ============================================================================
% CONFIDENCE-BASED EVALUATION PROMPT
% ============================================================================

\begin{figure}[htbp]
\centering
\begin{promptbox}[title={Confidence-Based Evaluation}]
Rate your confidence that the following statement is correct.

Provide a single number between 0.0 and 1.0:
- 0.0 means you are certain the statement is incorrect
- 0.5 means you are completely uncertain
- 1.0 means you are certain the statement is correct

For example:
- 0.1 = very likely incorrect
- 0.3 = probably incorrect  
- 0.5 = uncertain
- 0.7 = probably correct
- 0.9 = very likely correct

Respond with ONLY a decimal number between 0.0 and 1.0.
Do not include any explanation or additional text.

Statement: {assertion}

Confidence:
\end{promptbox}
\caption{Confidence-Based Evaluation Prompt. This approach elicits a numerical confidence rating which is then compared against different thresholds (0.5, 0.7, 0.9) to determine classification and coverage trade-offs.}
\label{fig:confidence-prompt}
\end{figure}

% ============================================================================
% EPISTEMIC POLICY PROJECTIONS
% ============================================================================

\begin{figure}[htbp]
\centering
\begin{tcolorbox}[colback=white,colframe=promptborder,title={Epistemic Policy Projections from Bilateral Truth Values}]
\small
\textbf{Classical Policy:} Abstains on both contradictions ⟨t,t⟩ and knowledge gaps ⟨f,f⟩\\[0.5em]
\begin{tabular}{l|l}
\hline
Bilateral Value & Classical Projection \\
\hline
⟨t,f⟩ & CORRECT \\
⟨f,t⟩ & INCORRECT \\
⟨t,t⟩ & ABSTAIN (contradiction) \\
⟨f,f⟩ & ABSTAIN (knowledge gap) \\
⟨t,e⟩, ⟨e,t⟩, ⟨f,e⟩, ⟨e,f⟩, ⟨e,e⟩ & ABSTAIN (epistemic failure) \\
\hline
\end{tabular}

\vspace{1em}
\textbf{Paracomplete Policy:} Abstains only on knowledge gaps ⟨f,f⟩\\[0.5em]
\begin{tabular}{l|l}
\hline
Bilateral Value & Paracomplete Projection \\
\hline
⟨t,f⟩ & CORRECT \\
⟨f,t⟩ & INCORRECT \\
⟨t,t⟩ & CORRECT (accepts contradiction as true) \\
⟨f,f⟩ & ABSTAIN (knowledge gap) \\
⟨t,e⟩, ⟨e,t⟩, ⟨f,e⟩, ⟨e,f⟩, ⟨e,e⟩ & ABSTAIN (epistemic failure) \\
\hline
\end{tabular}

\vspace{1em}
\textbf{Paraconsistent Policy:} Abstains only on contradictions ⟨t,t⟩\\[0.5em]
\begin{tabular}{l|l}
\hline
Bilateral Value & Paraconsistent Projection \\
\hline
⟨t,f⟩ & CORRECT \\
⟨f,t⟩ & INCORRECT \\
⟨t,t⟩ & ABSTAIN (contradiction) \\
⟨f,f⟩ & INCORRECT (treats knowledge gap as false) \\
⟨t,e⟩, ⟨e,t⟩, ⟨f,e⟩, ⟨e,f⟩, ⟨e,e⟩ & ABSTAIN (epistemic failure) \\
\hline
\end{tabular}
\end{tcolorbox}
\caption{Epistemic policy projections showing how bilateral truth values ⟨u,v⟩ map to final judgments under different epistemic frameworks. The Classical policy is most conservative, while Paracomplete and Paraconsistent policies offer different trade-offs between coverage and accuracy.}
\label{fig:epistemic-policies}
\end{figure}

% ============================================================================
% SUMMARY COMPARISON TABLE
% ============================================================================

\begin{figure}[htbp]
\centering
\begin{tcolorbox}[colback=white,colframe=promptborder,title={Comparison of Evaluation Approaches}]
\small
\begin{tabular}{l|p{3cm}|p{3cm}|p{3cm}|p{2.5cm}}
\hline
\textbf{Approach} & \textbf{Prompt Structure} & \textbf{Response Format} & \textbf{Coverage} & \textbf{Key Feature} \\
\hline
Bilateral & Two separate prompts (verification \& refutation) & VERIFIED/CANNOT VERIFY + REFUTED/CANNOT REFUTE & Variable (56\% avg) & Rich truth values ⟨u,v⟩ \\[0.5em]
\hline
Unilateral-Binary & Single direct question & CORRECT/INCORRECT & 100\% & Forced choice \\[0.5em]
\hline
Unilateral-Ternary & Single question with uncertainty option & CORRECT/INCORRECT/UNCERTAIN & Variable (52.5\% avg) & Explicit uncertainty \\[0.5em]
\hline
Confidence-0.5 & Single confidence rating & 0.0-1.0 (threshold at 0.5) & 100\% & Continuous scale \\[0.5em]
\hline
Confidence-0.7 & Single confidence rating & 0.0-1.0 (threshold at 0.7) & Variable (64.2\% avg) & Higher threshold \\[0.5em]
\hline
Confidence-0.9 & Single confidence rating & 0.0-1.0 (threshold at 0.9) & Variable (46.7\% avg) & Highest threshold \\[0.5em]
\hline
\end{tabular}
\end{tcolorbox}
\caption{Summary comparison of all evaluation approaches showing their prompt structure, response formats, coverage characteristics, and distinguishing features.}
\label{fig:approach-comparison}
\end{figure}

\end{document}